\chapter{Faire des mathématiques en conversation : \LaTeX{} et Lean}

La combinaison la plus utile du projet a été celle-ci : discuter d'un énoncé avec
l'assistant, le faire vérifier par Lean, puis le rédiger en \LaTeX{} à partir de la preuve
vérifiée. Chacun des trois corrige les deux autres. La conversation propose, Lean tranche,
et la rédaction oblige à comprendre ce qui a été démontré.

\section{La boucle}

\begin{enumerate}
  \item \textbf{Poser la question en mathématicien.} On écrit l'énoncé dans la
        conversation, avec ses hypothèses, comme à un collègue. L'assistant cherche une
        démonstration, propose des cas particuliers, signale ce qui manque.
  \item \textbf{Tester avant de prouver.} Quand l'énoncé porte sur des objets finis,
        un script vérifie d'abord les petits cas. Un contre-exemple trouvé en dix
        secondes épargne une journée de preuve.
  \item \textbf{Énoncer en Lean.} L'assistant écrit l'énoncé dans Lean, avec
        \texttt{sorry} comme preuve, et le fait compiler. C'est le moment où les
        hypothèses implicites apparaissent : Lean refuse un énoncé qui ne dit pas sur
        quel anneau, quelle catégorie, quel ensemble il porte.
  \item \textbf{Prouver.} L'assistant comble le \texttt{sorry}, en s'aidant des outils de
        recherche de la bibliothèque, compile, lit les erreurs, recommence. On vérifie à
        la fin les axiomes (\cmd{\#print axioms}).
  \item \textbf{Rédiger.} À partir de la preuve qui compile, l'assistant écrit la
        démonstration en \LaTeX, dans la langue de l'article, en suivant les étapes de la
        preuve formelle. Le texte renvoie au fichier et au nom du théorème.
\end{enumerate}

\begin{regle}
Le fichier Lean fait foi. Si la rédaction et la preuve divergent, c'est la rédaction qu'on
corrige ; si l'énoncé de l'article ne correspond pas à celui de Lean, on le dit dans
l'article.
\end{regle}

\section{Un exemple}

Une liste de problèmes ouverts contenait la question suivante : un arrangement de
pseudo-droites du plan projectif, dont chaque pseudo-droite porte au moins trois sommets,
est-il déterminé par son graphe étiqueté ? Une formalisation antérieure réglait le cas où
chaque triangle a une quatrième droite évitant ses sommets ; restait le cas des trois
faisceaux.

\begin{exemple}
En conversation, l'assistant a reformulé la question en termes de chirotopes : le graphe
fixe, le long de chaque droite, l'ordre cyclique des sommets, donc les signes à un
retournement près. Il en a tiré une fonction de signes qui doit valoir~$1$. Un script a
vérifié la version linéaire sur $\mathbb F_2$ pour tous les arrangements d'au plus neuf
droites (37\,213 treillis, aucune obstruction). Le cas difficile se réglait en coupant le
triangle en deux le long d'une droite bien choisie. Lean a ensuite démontré l'algèbre des
signes, puis le pas de rang~$2$. L'article, enfin, donne la démonstration en une page, avec
une figure, et dit ce qui reste hors de Lean : le théorème de représentation
topologique.
\end{exemple}

Deux remarques sur cet exemple. La première idée venait de la conversation, pas de Lean ;
Lean n'a fait que la rendre sûre. Et l'histoire du travail se lit dans le dépôt : le script
de vérification, le fichier Lean, le commit qui modifie l'article, chacun avec son message.

\section{Ce que Lean ne dit pas}

Une preuve formelle établit la lecture qu'on en a faite. Elle ne dit pas que cette lecture
est celle de l'auteur d'un texte qu'on édite, ni que le résultat est nouveau
(chapitre~\ref{ch:litterature}). Elle ne dit pas non plus que la rédaction en \LaTeX{} est
claire. Un lecteur humain reste nécessaire pour cela, et la passe d'humanisation du
chapitre précédent ne le remplace pas.

\begin{warningblock}
Un assistant peut \og démontrer \fg{} un énoncé en Lean en le rendant trivial : une
hypothèse trop forte, une définition vide, un \texttt{sorry} caché dans une définition. On
lit l'énoncé du théorème avant de lire sa preuve, et l'on vérifie qu'il dit bien ce que
l'article affirme.
\end{warningblock}
